\documentclass[10pt,a4paper]{amsart}
\usepackage[margin=19mm]{geometry}
\usepackage{amsmath,amssymb,mathtools}
\usepackage[scr=rsfs,cal=euler]{mathalfa}
\usepackage{xfrac}
\usepackage[unicode,hidelinks,bookmarks=false]{hyperref}
\newcommand{\ie}{i.e.\ }
\newcommand{\cf}{cf.\ }
\newcommand{\resp}{respectively\ }
\newcommand{\Subsection}{\subsection}
\providecommand{\nobreakdash}{\nobreak\hbox{-}}
\setlength{\emergencystretch}{2em}
\multlinegap0pt
\usepackage{amssymb}
\newtheorem{theorem}{Theorem}
\title{An observation on factorizations of finite groups}
\author{George M. Bergman}
\date{Source: arXiv:2608.08928v2 (CC BY 4.0)}
\begin{document}
\maketitle

\noindent\emph{Source and licence.} An observation on factorizations of finite groups. Version arXiv:2608.08928v2 (CC BY 4.0), distributed by arXiv under the Creative Commons Attribution 4.0 International licence, http://creativecommons.org/licenses/by/4.0/, as recorded in the arXiv licence metadata for this exact version and shown on the abstract page https://arxiv.org/abs/2608.08928v2. Full licence notice: \texttt{licenses/arxiv-2608.08928-CC-BY-4.0.txt}.\par\smallskip

\noindent\textit{Editorial scope.} Counted unit: the article's theorem T.main (source0.tex lines 126--148). The article's complete original proof of it is delivered with it (source0.tex lines 150--204, one \texttt{proof} environment of 55 lines). No mathematics is written, repaired or re-derived: the statement and the proof are the article's own lines, in the article's own notation, and no retained line is substituted or re-flowed.  No further source-local context is needed: every cross-reference of every retained line resolves inside the two delivered spans.  Nothing was omitted from any retained span, so the package carries no omission record.  No retained line contains a citation, so the wrapper carries no bibliography and no bibliography entry.  No retained line contains a figure environment, an image-inclusion command or drawing code, so no image is delivered and the package declares no asset dependency.  Editorial formatting only, in addition: an AMS \texttt{amsart} preamble that restates the article's own declarations and macros used by the retained lines, namely the environments \texttt{theorem} on separate counters, together with the standard packages those lines need.  The retained environments print in this document as Theorem 1 in order of appearance, whereas the article prints the same items with the numbering of its own full document.  The complete native source of the article as distributed in the arXiv e-print for this exact version is bundled unchanged as \texttt{sources/207202/source0.tex}.
\begin{theorem}\label{T.main}
Let $G$ be a finite group, and
%
\begin{equation}\begin{minipage}[c]{35pc}\label{d.G0Gn}
$\{e\}\,=\,G_0\ <\ G_1\ <\ \dots\ <\ G_n\,=\,G$ \ $(n\geq 0)$
\end{minipage}\end{equation}
%
a chain of subgroups of $G.$

Then for any factorization $|G| = a_1\,a_2$ with $a_1$ and $a_2$
positive integers, a {\em sufficient}
condition for there to exist subsets $A_1,\,A_2\subseteq G$
having cardinalities $|A_1|=a_1,$ $|A_2|=a_2$ and satisfying
%
\begin{equation}\begin{minipage}[c]{35pc}\label{d.A1A2}
$G\ =\ A_1\,A_2$
\end{minipage}\end{equation}
%
is that there exist a subset $s_1\subseteq\{1,\dots,n\}$
such that $\prod_{i\in s_1}\,|G_i:\,G_{i-1}| = a_1,$
equivalently, such that, letting $s_2=\{1,\dots,n\}\setminus s_1,$
we have $\prod_{i\in s_2}\,|G_i:\,G_{i-1}| = a_2.$
\end{theorem}

\begin{proof}
Given a finite group $G,$ a chain~\eqref{d.G0Gn}, and
complementary sets $s_1,\ s_2\subseteq\{1,\dots,n\},$
we shall show by induction on $j$ that for each $j\leq n,$
%
\begin{equation}\begin{minipage}[c]{35pc}\label{d.A1jA2j}
there exists a factorization $G_j=A_{1,j}\ A_{2,j}$ such that \\
$|A_{1,j}|\ =\ \prod_{i\in s_1\cap\{1,\dots,j\}}\,|G_i:\,G_{i-1}|$ \ and
\ $|A_{2,j}|\ =\ \prod_{i\in s_2\cap\{1,\dots,j\}}\,|G_i:\,G_{i-1}|.$
\end{minipage}\end{equation}

This is trivially true for $j=0.$
(In that case the numerical products shown are over the empty set,
hence their values are both $1,$ and
$G_0=\{e\}$ indeed has the factorization $\{e\}=\{e\}\{e\}=
(\prod_\emptyset)\,(\prod_\emptyset).)$

So let $0\leq j<n,$ and assume inductively that
we have a factorization $G_j=A_{1,j}\ A_{2,j}$ as in~\eqref{d.A1jA2j}.

The integer $j\,{+}\,1$ occurs in exactly
one of the subsets $s_1,$ $s_2.$
Without loss of generality let us assume
%
\begin{equation}\begin{minipage}[c]{35pc}\label{d.i_in_s1}
$j\,{+}\,1\,\in\,s_1.$
\end{minipage}\end{equation}
%
(``Without loss of generality'' because our inductive
hypothesis and desired conclusion are symmetric
with respect to reversal of order of multiplication, which
interchanges the roles of $s_1$ and $s_2.)$
Assuming~\eqref{d.i_in_s1},
%
\begin{equation}\begin{minipage}[c]{35pc}\label{d.Cj+1}
let $C_{j+1}$ be any set of representatives
of the left cosets of $G_j$ in $G_{j+1},$ and let \\
$A_{1,j+1}\ =\ C_{j+1}\,A_{1,j},$ \ $A_{2,j+1}\ =\ A_{2,j}.$
\end{minipage}\end{equation}

Given that $A_{1,j}$ and $A_{2,j}$
have the cardinalities indicated in~\eqref{d.A1jA2j},
and assuming~\eqref{d.i_in_s1}, we see that
the new sets $A_{1,j+1}$ and $A_{2,j+1}$ will
have cardinalities as in the $j\,{+}\,1$ case of~\eqref{d.A1jA2j}.
Moreover, in view of~\eqref{d.Cj+1},
%
\begin{equation}\begin{minipage}[c]{35pc}\label{d.CiA1}
$A_{1,j+1}\,A_{2,j+1}\ =\ C_{j+1}\ A_{1,j}\ A_{2,j}\ =\ C_{j+1}\ G_j\
=\ G_{j+1}.$
\end{minipage}\end{equation}

So we have proved the $j\,{+}\,1$ case of~\eqref{d.A1jA2j}; so by
induction,~\eqref{d.A1jA2j} holds for $j=n,$ our desired conclusion.
\end{proof}

\end{document}
